\begin{lemma}
\label{lemma-slice-given-element}
Let $f : X \to S$ be a morphism of schemes.
Let $x \in X$ be a point with image $s \in S$.
Let $h \in \mathfrak m_x \subset \mathcal{O}_{X, x}$.
Assume
\begin{enumerate}
\item $f$ is locally of finite presentation,
\item $f$ is flat at $x$, and
\item the image $\overline{h}$ of $h$ in
$\mathcal{O}_{X_s, x} = \mathcal{O}_{X, x}/\mathfrak m_s\mathcal{O}_{X, x}$
is a nonzerodivisor.
\end{enumerate}
Then there exists an affine open neighbourhood $U \subset X$ of $x$
such that $h$ comes from $h \in \Gamma(U, \mathcal{O}_U)$ and such
that $D = V(h)$ is an effective Cartier divisor in $U$ with $x \in D$ and
$D \to S$ flat and locally of finite presentation.
\end{lemma}

\begin{proof}
We are going to prove this by reducing to the Noetherian case.
By openness of flatness (see
Theorem \ref{theorem-openness-flatness})
we may assume, after replacing $X$ by an
open neighbourhood of $x$, that $X \to S$ is flat.
We may also assume that $X$ and $S$ are affine.
After possible shrinking $X$ a bit we may assume that there exists
an $h \in \Gamma(X, \mathcal{O}_X)$ which maps to our given $h$.

\medskip\noindent
We may write $S = \Spec(A)$ and we may write $A = \colim_i A_i$
as a directed colimit of finite type $\mathbf{Z}$ algebras.
Then by
Algebra, Lemma \ref{algebra-lemma-flat-finite-presentation-limit-flat}
or
Limits, Lemmas \ref{limits-lemma-descend-finite-presentation},
\ref{limits-lemma-descend-affine-finite-presentation}, and
\ref{limits-lemma-descend-finite-presentation}
we can find a cartesian diagram
$$
\xymatrix{
X \ar[r] \ar[d]_f & X_0 \ar[d]^{f_0} \\
S \ar[r] & S_0
}
$$
with $f_0$ flat and of finite presentation, $X_0$ affine, and
$S_0$ affine and Noetherian. Let $x_0 \in X_0$, resp.\ $s_0 \in S_0$
be the image of $x$, resp.\ $s$. We may also assume there exists an element
$h_0 \in \Gamma(X_0, \mathcal{O}_{X_0})$ which restricts to $h$ on $X$.
(If you used the algebra reference above then this is clear; if you used
the references to the chapter on limits then this follows from
Limits, Lemma \ref{limits-lemma-descend-finite-presentation}
by thinking of $h$ as a morphism $X \to \mathbf{A}^1_S$.)
Note that $\mathcal{O}_{X_s, x}$ is a localization of
$\mathcal{O}_{(X_0)_{s_0}, x_0} \otimes_{\kappa(s_0)} \kappa(s)$, so that
$\mathcal{O}_{(X_0)_{s_0}, x_0} \to \mathcal{O}_{X_s, x}$ is a flat
local ring map, in particular faithfully flat. Hence the image
$\overline{h}_0 \in \mathcal{O}_{(X_0)_{s_0}, x_0}$
is contained in $\mathfrak m_{(X_0)_{s_0}, x_0}$ and is a nonzerodivisor.
We claim that after replacing $X_0$ by a principal open neighbourhood of
$x_0$ the element $h_0$ is a nonzerodivisor in
$B_0 = \Gamma(X_0, \mathcal{O}_{X_0})$ such that $B_0/h_0B_0$ is flat
over $A_0 = \Gamma(S_0, \mathcal{O}_{S_0})$.
If so then
$$
0 \to B_0 \xrightarrow{h_0} B_0 \to B_0/h_0B_0 \to 0
$$
is a short exact sequence of flat $A_0$-modules. Hence this remains exact
on tensoring with $A$ (by
Algebra, Lemma \ref{algebra-lemma-flat-tor-zero})
and the lemma follows.

\medskip\noindent
It remains to prove the claim above. The corresponding algebra statement
is the following (we drop the subscript ${}_0$ here):
Let $A \to B$ be a flat, finite type ring map of Noetherian rings.
Let $\mathfrak q \subset B$ be a prime lying over $\mathfrak p \subset A$.
Assume $h \in \mathfrak q$ maps to a nonzerodivisor in
$B_{\mathfrak q}/\mathfrak p B_{\mathfrak q}$.
Goal: show that after possible replacing $B$ by $B_g$ for some
$g \in B$, $g \not \in \mathfrak q$ the element $h$ becomes a nonzerodivisor
and $B/hB$ becomes flat over $A$. By
Algebra, Lemma \ref{algebra-lemma-grothendieck}
we see that $h$ is a nonzerodivisor in $B_{\mathfrak q}$ and that
$B_{\mathfrak q}/hB_{\mathfrak q}$ is flat over $A$.
By openness of flatness, see
Algebra, Theorem \ref{algebra-theorem-openness-flatness}
or
Theorem \ref{theorem-openness-flatness}
we see that $B/hB$ is flat over $A$ after replacing $B$ by $B_g$ for some
$g \in B$, $g \not \in \mathfrak q$. Finally, let
$I = \{b \in B \mid hb = 0\}$ be the annihilator of $h$. Then
$IB_{\mathfrak q} = 0$ as $h$ is a nonzerodivisor in $B_{\mathfrak q}$.
Also $I$ is finitely generated as $B$ is Noetherian.
Hence there exists a $g \in B$, $g \not \in \mathfrak q$ such that $IB_g = 0$.
After replacing $B$ by $B_g$ we see that $h$ is a nonzerodivisor.
\end{proof}